% Generated by benchmarks/compile_sw07_profile_comparison.py. Do not edit.
\begin{table}[t]
\centering
\caption{Matched SW07 CPU--GPU likelihood and gradient profile}
\label{tab:sw07-cpu-gpu-profile}
\small
\begin{tabular}{lrrr}
\toprule
Metric (seconds) & CPU & GPU & CPU/GPU \\
\midrule
Likelihood first call & 7.433 & 44.907 & 0.17 \\
Likelihood steady median & 0.015 & 0.065 & 0.22 \\
Gradient first call & 330.261 & 290.327 & 1.14 \\
Gradient steady median & 0.080 & 1.472 & 0.05 \\
\bottomrule
\end{tabular}

\vspace{0.4em}
\parbox{0.96\textwidth}{\footnotesize\textit{Notes:} Both backends evaluate the same float64, 36-parameter likelihood on the 180-quarter sample, solve the steady state at every draw, and use theoretical Kalman initialization. First-call times include JAX/XLA compilation; steady medians exclude compilation. CPU/GPU is the ratio of elapsed times, so values above one favor the GPU. The likelihood discrepancy is $3.87\times 10^{-12}$; the maximum gradient-component discrepancy is $5.43\times 10^{-10}$. The CPU environment is macOS arm64 with JAX 0.4.30 and NumPyro 0.19.0; the GPU environment is Linux x86\_64 on an RTX 5090 D with JAX 0.11.1 and NumPyro 0.21.0. The timing ratio is therefore an observational deployment comparison, not a controlled hardware-only ablation. These microprofiles do not measure effective samples per second.}
\end{table}
